\documentclass[10pt,a4paper]{amsart}
\usepackage[margin=19mm]{geometry}
\usepackage{amsmath,amssymb,mathtools}
\usepackage[scr=rsfs,cal=euler]{mathalfa}
\usepackage{xfrac}
\usepackage[unicode,hidelinks,bookmarks=false]{hyperref}
\newcommand{\ie}{i.e.\ }
\newcommand{\cf}{cf.\ }
\newcommand{\resp}{respectively\ }
\newcommand{\Subsection}{\subsection}
\providecommand{\nobreakdash}{\nobreak\hbox{-}}
\setlength{\emergencystretch}{2em}
\multlinegap0pt
\usepackage[shortlabels]{enumitem}
\usepackage{amssymb}
\usepackage{mathtools}
\usepackage{bm}
\usepackage{xcolor}
\usepackage{tikz}
\usepackage{braket}
\newtheorem{prop}{Proposition}
\newcommand{\Lk}{\mathrm{Lk}}
\title{An exact spacetime polymer gas for finite-temperature $\mathbb Z_N$ homological quantum code}
\author{Nafiz Ishtiaque, Shanto Chakroborty}
\date{Source: arXiv:2605.09122v1 (CC BY 4.0)}
\begin{document}
\maketitle

\noindent\emph{Source and licence.} An exact spacetime polymer gas for finite-temperature $\mathbb Z_N$ homological quantum code. Version arXiv:2605.09122v1 (CC BY 4.0), distributed by arXiv under the Creative Commons Attribution 4.0 International licence, http://creativecommons.org/licenses/by/4.0/, as recorded in the arXiv licence metadata for this exact version and shown on the abstract page https://arxiv.org/abs/2605.09122v1. Full licence notice: \texttt{licenses/arxiv-2605.09122-CC-BY-4.0.txt}.\par\smallskip

\noindent\textit{Editorial scope.} Counted unit: the article's prop prop:choice-dependence (source0.tex lines 4891--4938). The article's complete original proof of it is delivered with it (source0.tex lines 4940--5029, one \texttt{proof} environment of 90 lines). No mathematics is written, repaired or re-derived: the statement and the proof are the article's own lines, in the article's own notation, and no retained line is substituted or re-flowed.  Retained uncounted as the source-local context the proof needs: lines 4668--4673; lines 4729--4734; lines 4736--4746; lines 4748--4751; lines 4781--4790; lines 4846--4853.  Nothing was omitted from any retained span, so the package carries no omission record.  No retained line contains a citation, so the wrapper carries no bibliography and no bibliography entry.  No retained line contains a figure environment, an image-inclusion command or drawing code, so no image is delivered and the package declares no asset dependency.  Editorial formatting only, in addition: an AMS \texttt{amsart} preamble that restates the article's own declarations and macros used by the retained lines, namely the environments \texttt{prop} on separate counters, together with the standard packages those lines need.  The retained environments print in this document as Proposition 1 in order of appearance, whereas the article prints the same items with the numbering of its own full document.  The complete native source of the article as distributed in the arXiv e-print for this exact version is bundled unchanged as \texttt{sources/200197/source0.tex}.
\begin{equation}
	I_X^{(r+1)}(\partial \Xi^\vee,\Sigma)
	=
	I_X^{(r)}(\Xi^\vee,\partial \Sigma).
	\label{eq:intersection-Stokes}
\end{equation}

	\begin{equation}
		s_p:H_p(X)\rightarrow Z_p(X),
		\qquad
		s^\vee_{n-p-1}:H_{n-p-1}(X^\vee)\rightarrow Z_{n-p-1}(X^\vee),
		\label{eq:sections}
	\end{equation}

	\begin{align}
		K_p:B_p(X)\rightarrow&\; \mathcal C_{p+1}(X),
		&\qquad
		\partial K_p=&\;\mathrm{id}_{B_p(X)},
		\label{eq:Kp}
	\\
		K^\vee_{n-p-1}:B_{n-p-1}(X^\vee)\rightarrow&\; \mathcal C_{n-p}(X^\vee),
		&
		\partial K^\vee_{n-p-1}=&\;\mathrm{id}_{B_{n-p-1}(X^\vee)},
		\label{eq:Kdual}
	\end{align}

	\begin{equation}
		B_H:H_{n-p-1}(X^\vee)\times H_p(X)\rightarrow \mathbb Z_N.
		\label{eq:BH}
	\end{equation}

\begin{equation}
	\Lk_X(\mu^\vee,\nu)
	:=
	I_X^{(p+1)}\left(\mu^\vee,\,K_p(\nu_{\mathrm ex})\right)
	+
	I_X^{(p)}\left(K^\vee_{n-p-1}(\mu^\vee_{\mathrm ex}),\,\nu_{\mathrm h}\right)
	+
	B_H\left([\mu^\vee],[\nu]\right).
	\label{eq:generalized-linking}
\end{equation}

\begin{equation}
	\Lk_X(\mu^\vee,\nu)
	=
	I_X^{(p+1)}\left(\mu^\vee,K_p(\nu)\right)
	=
	I_X^{(p)}\left(K^\vee_{n-p-1}(\mu^\vee),\nu\right),
	\label{eq:boundary-linking}
\end{equation}

\begin{prop}
	\label{prop:choice-dependence}
	Fix a single choice of the filling operators \eqref{eq:Kp} and \eqref{eq:Kdual}. Let
	\begin{equation}
		(s_p,s^\vee_{n-p-1},B_H)
		\qquad\text{and}\qquad
		(\widetilde s_p,\widetilde s^\vee_{n-p-1},\widetilde B_H)
	\end{equation}
	be two choices of sections \eqref{eq:sections} and homological correction term
	\eqref{eq:BH}, and let \(\Lk_X\) and \(\widetilde{\Lk}_X\) be the corresponding
	generalized linking pairings \eqref{eq:generalized-linking}. Then:
	\begin{enumerate}[label=(\roman*)]
		\item\label{prop:choice-dependence-item1}
		There exists a bilinear map
		\begin{equation}
			C_H:H_{n-p-1}(X^\vee)\times H_p(X)\to \mathbb Z_N
			\label{eq:DeltaH}
		\end{equation}
		such that
		\begin{equation}
			\widetilde{\Lk}_X(\mu^\vee,\nu)-\Lk_X(\mu^\vee,\nu)
			=
			C_H\left([\mu^\vee],[\nu]\right)
			\label{eq:choice-dependence}
		\end{equation}
		for all
		\(\mu^\vee\in Z_{n-p-1}(X^\vee)\) and \(\nu\in Z_p(X)\). In particular, the two
		pairings agree whenever either \(\mu^\vee\) or \(\nu\) is a boundary.
		
		\item\label{propItem:bdryLink}
		If \(\mu^\vee\in B_{n-p-1}(X^\vee)\) and \(\nu\in B_p(X)\), then
		\begin{equation}
			\Lk_X(\mu^\vee,\nu)
			=
			I_X^{(p+1)}\left(\mu^\vee,\Sigma\right)
			=
			I_X^{(p)}\left(T^\vee,\nu\right),
			\qquad
			\partial\Sigma=\nu,\ \partial T^\vee=\mu^\vee.
			\label{eq:boundary-linking-canonical-merged}
		\end{equation}
		This value is independent of all auxiliary choices used in the definition of
		\(\Lk_X\), including the sections, the homological correction term, and the
		filling chains \(\Sigma\), \(T^\vee\). In particular, the linking number of two
		boundaries is canonical and does not require the assumption that \(N\) be
		prime.
	\end{enumerate}
\end{prop}

\begin{proof}
	Set
	\begin{equation}
		\begin{aligned}
			\alpha&:=\widetilde s_p-s_p:H_p(X)\to Z_p(X),
			\\
			\beta&:=\widetilde s^\vee_{n-p-1}-s^\vee_{n-p-1}:H_{n-p-1}(X^\vee)\to Z_{n-p-1}(X^\vee),
		\end{aligned}
		\label{eq:secDiffs}
	\end{equation}
	and
	\begin{equation}
		C_H^{0}:=\widetilde B_H-B_H:
		H_{n-p-1}(X^\vee)\times H_p(X)\to \mathbb Z_N.
	\end{equation}
	Since both \(s_p\) and \(\widetilde s_p\) are sections of
	\(Z_p(X)\twoheadrightarrow H_p(X)\), the image of \(\alpha\) lies in \(B_p(X)\).
	Likewise, the image of \(\beta\) lies in \(B_{n-p-1}(X^\vee)\).
	
	Writing the homological and exact parts defined by the two choices of sections, we have
	\begin{equation}
		\begin{aligned}
			\widetilde\nu_{\mathrm h}&=\nu_{\mathrm h}+\alpha([\nu]),
			&\qquad
			\widetilde\nu_{\mathrm ex}&=\nu_{\mathrm ex}-\alpha([\nu]),
			\\
			\widetilde\mu^\vee_{\mathrm h}&=\mu^\vee_{\mathrm h}+\beta([\mu^\vee]),
			&
			\widetilde\mu^\vee_{\mathrm ex}&=\mu^\vee_{\mathrm ex}-\beta([\mu^\vee]).
		\end{aligned}
	\end{equation}
	Substituting these into the definition of \(\widetilde{\Lk}_X\), expanding, and using
	\eqref{eq:intersection-Stokes} to move fillings of exact terms from one side to the
	other, we obtain
	\begin{equation}
		\begin{aligned}
			\widetilde{\Lk}_X(\mu^\vee,\nu)-\Lk_X(\mu^\vee,\nu)
			=&\;
			-I_X^{(p+1)}\left(\mu^\vee_{\mathrm h},K_p(\alpha([\nu]))\right)
			-I_X^{(p)}\left(K^\vee_{n-p-1}(\beta([\mu^\vee])),\nu_{\mathrm h}\right)
			\\
			&\;
			-I_X^{(p)}\left(K^\vee_{n-p-1}(\beta([\mu^\vee])),\alpha([\nu])\right)
			+C_H^{0}([\mu^\vee],[\nu]).
		\end{aligned}
		\label{eq:choice-proof}
	\end{equation}
	Now \(\mu^\vee_{\mathrm h}=s^\vee_{n-p-1}([\mu^\vee])\) and
	\(\nu_{\mathrm h}=s_p([\nu])\), so the right-hand side depends only on
	\([\mu^\vee]\) and \([\nu]\), and it is bilinear because all maps appearing in
	\eqref{eq:choice-proof} are linear or bilinear. This proves
	\eqref{eq:choice-dependence} and therefore
	\ref{prop:choice-dependence-item1}. If either \([\mu^\vee]=0\) or \([\nu]=0\),
	then all terms on the right-hand side of \eqref{eq:choice-proof} vanish, so the two
	pairings agree whenever either argument is a boundary.
	
	For \ref{propItem:bdryLink}, if
	\(\mu^\vee\in B_{n-p-1}(X^\vee)\) and \(\nu\in B_p(X)\), then
	\eqref{eq:boundary-linking} gives
	\begin{equation}
		\Lk_X(\mu^\vee,\nu)
		=
		I_X^{(p+1)}\left(\mu^\vee,\Sigma\right)
		=
		I_X^{(p)}\left(T^\vee,\nu\right),
		\qquad
		\partial\Sigma=\nu,\ \partial T^\vee=\mu^\vee.
	\end{equation}
	By \ref{prop:choice-dependence-item1}, this value is independent of the chosen
	sections and homological correction term, because both homology classes vanish.
	It remains only to check independence of the fillings. If
	\(\Sigma,\Sigma'\in\mathcal C_{p+1}(X)\) both satisfy
	\(\partial\Sigma=\partial\Sigma'=\nu\), then \(\Sigma-\Sigma'\in Z_{p+1}(X)\), so
	using \(\mu^\vee=\partial T^\vee\) and \eqref{eq:intersection-Stokes},
	\begin{equation}
		I_X^{(p+1)}\left(\mu^\vee,\Sigma-\Sigma'\right)
		=
		I_X^{(p+1)}\left(\partial T^\vee,\Sigma-\Sigma'\right)
		=
		I_X^{(p)}\left(T^\vee,\partial(\Sigma-\Sigma')\right)
		=
		0.
	\end{equation}
	Hence \(I_X^{(p+1)}(\mu^\vee,\Sigma)\) is independent of \(\Sigma\). The argument for
	\(I_X^{(p)}(T^\vee,\nu)\) is identical. Therefore the common value
	\eqref{eq:boundary-linking-canonical-merged} is canonical. Since its definition uses
	only ordinary filling-intersection of two boundaries, it does not depend on the
	prime-\(N\) hypothesis introduced to construct a generalized pairing on arbitrary
	cycles.
\end{proof}

\end{document}
